\chapter{多项式映射}\label{chap:polynomial-maps}

\noindent\textbf{第 \thechapter{} 章学习目标}
\begin{itemize}
  \item 定义仿射簇之间的多项式映射并举例, 包括同构的例子.
  \item 定义与多项式映射关联的拉回同态, 并在具体例子中计算它.
  \item 利用拉回的双射性证明: 仿射簇同构当且仅当它们的坐标环同构.
  \item 描述代数与几何的等价, 以及它怎样在 $K$-代数和仿射簇的内在性质之间进行转换.
\end{itemize}

上一章给出了一个对应
\[
  \begin{aligned}
    \{\text{仿射簇}\}&\to\{\text{有限生成的既约 }K\text{-代数}\}\\
    X&\mapsto K[X]
  \end{aligned}
\]
并确认了它是满射. 但它是单射吗? 也就是说, 如果 $K[X]=K[Y]$, 是否必有 $X=Y$?

这个问题比乍看起来更微妙. 要回答它, 必须准确决定何时两个 $K$-代数是\textquotedbl{}同一个\textquotedbl{}: 例如, $K[x]$ 与 $K[y]$ 是同一个吗? 从字面上说, 不是. 但我们是如此习惯于把同构的环看成相同的, 所以读者觉得这个答案不能令人满意, 也是可以理解的. 为抓住同构就是\textquotedbl{}相同\textquotedbl{}这一直觉, 可以重新表述问题: 如果 $K[X]\cong K[Y]$, 是否必有 $X\cong Y$?

要处理这个问题, 必须首先精确定义两个仿射簇\textquotedbl{}同构\textquotedbl{}是什么意思, 从仿射簇之间保持结构的\textquotedbl{}态射\textquotedbl{}这一概念着手. 每当引入一种新类型的数学对象 (例如群, 环, 拓扑空间) 时, 都应做这件事: 问这些对象之间哪些映射保持有关的结构. 对群来说, 有关的映射是群同态; 对环来说, 是环同态; 对拓扑空间来说, 是连续映射. 因为代数几何研究多项式, 仿射簇之间有关的映射是多项式映射, 或许不足为奇. 我们将在本章引入这些映射.

有了多项式映射的定义, 就能说清仿射簇 $X$ 与 $Y$ 同构是什么意思, 并能证明 $X\cong Y$ 当且仅当 $K[X]\cong K[Y]$. 这一陈述称为\textquotedbl{}代数与几何的等价\textquotedbl{}, 是本章的目标, 也是代数与几何之间的词典的核心.

\par\smallskip
{\footnotesize
\noindent\textcopyright{} The Author(s) 2025\par
\noindent E. Clader, D. Ross, \emph{Beginning in Algebraic Geometry}, Undergraduate Texts in Mathematics,
\href{https://doi.org/10.1007/978-3-031-88819-9_5}{\nolinkurl{https://doi.org/10.1007/978-3-031-88819-9_5}}.\par}
